\section{Supporting synthetic reduced-view study}
\label{sec:hlt-support}
\label{sec:supporting-results}

This appendix documents the exploratory experiments that preceded the
offline study. The broad relation screen characterizes performance under
one specified synthetic degradation and fixes the selected relation set;
the architecture comparisons provide supporting evidence about value
messages and bias sharing. They are not measurements of real trigger
performance or a substitute for the offline architectural ablations.

\subsection{Synthetic input definition and scope}
\label{sec:hlt-inputs}

The reduced view, called HLT-like for continuity with the original
campaign, is generated from the offline constituents using the
repository's \texttt{fixed\_hlt\_v1} profile at strength 0.6. It applies
constituent thresholding, local angular merging, kinematics- and
density-dependent inefficiency, momentum and angular smearing, and local
angular reassignment. The constituent $p_T$ threshold is
$0.78\,\mathrm{GeV}$ and the merging radius is 0.0081. One cached
realization per split, generated with a fixed random seed, is shared by all
configurations and is not resampled during training.

This is a controlled input perturbation, not a simulation of a particular
experiment's trigger reconstruction. Surviving particles retain their
supplied identification, charge, impact parameters, and uncertainties;
merged particles inherit the dominant constituent's nonkinematic entries.
Tracking and identification resolutions are not independently resimulated.

The study uses the development and earlier-test split sizes in
Table~\ref{tab:experimental-splits}, with the same underlying jet identities
as the original offline campaign. Particle features, axes, four-vectors,
and pair descriptors are computed from the reduced view; normalization is
fitted only on its training sample. No privileged offline features or
teacher predictions are supplied to the reduced-view classifier.
The backbone and training protocol follow Sections~\ref{sec:experimental-models}
and~\ref{sec:training-protocol}; the initial screen uses one common seed,
and architecture confirmation uses three matched seeds.

\subsection{Relation-family screening protocol}
\label{sec:hlt-screen}

The initial screen contains 21 configurations trained with one common random seed. All use
shared pair bias without EdgeValue. Three are controls: the standard
baseline; a widened standard-feature pair encoder matched as closely as
possible to the largest incremental relation-encoder parameter budget; and
a full-relation input-shape control whose additional relation channels are
zero. These controls examine additional capacity and the presence of the
expanded input structure, without supplying the corresponding relation
information.

The remaining 18 configurations comprise six individual relation families
(Table~\ref{tab:relation-families}) and twelve predetermined combinations.
The seven pairs are PT+TRACK, TRACK+PID, TRACK+CHARGE, PID+CHARGE,
PT+DENSITY, PT+REGION, and TRACK+REGION. Higher-order combinations are
TRACK+PID+CHARGE, PT+TRACK+DENSITY, PT+TRACK+REGION,
PT+TRACK+PID+CHARGE, and all six families together. Every such configuration
retains the standard four pair features.

\begin{table}[htbp]
\centering
\small
\begin{tabularx}{\linewidth}{@{}lXr@{}}
\toprule
Family & Additional particle-pair information & Channels \\
\midrule
PT & Endpoint momentum fractions and ranks, with directed differences and asymmetry. & 8 \\
TRACK & Impact parameters, uncertainties and significances, with pair compatibility and directed differences. & 12 \\
PID & Ordered pairs of particle-identification categories. & 8 \\
CHARGE & Ordered charge states, products, and comparison indicators. & 6 \\
DENSITY & Multiscale local multiplicity, momentum activity and composition around the endpoints. & 12 \\
REGION & Shared cluster membership, common-ancestor descriptors, and endpoint cluster properties at $K=2,4,8$. & 12 \\
\bottomrule
\end{tabularx}
\caption{Relation families considered in the HLT-like screen. Channel counts
refer to their encoded outputs before concatenation with the standard four
pair features, not their raw descriptor dimensions. REGION uses a
deterministic beam-free angular clustering tree, not generator ancestry.}
\label{tab:relation-families}
\end{table}

Screening uses the comparison-validation split, not the test set. Among
eligible scientific configurations, candidates within $10^{-4}$ absolute
accuracy (0.01 percentage points) of the highest validation accuracy are
ordered by lower cross-entropy, then fewer parameters, then run identifier.
The same rule generates the ranking used for follow-up selection. This
procedure selects PT+TRACK+REGION; this feature set is fixed before the
offline campaign. The resulting selected pair input has
$4+8+12+12=36$ channels. Its clustering descriptors are reconstructed from
the available input particles and do not supply truth decay or vertex
information.

\subsection{Multi-seed architecture and control protocol}
\label{sec:hlt-confirmation}

Follow-up comparisons use three matched training seeds. The central architectural
comparison is standard versus selected pair features, each with (i) shared
bias and no EdgeValue, (ii) layerwise bias and no EdgeValue, or (iii)
layerwise bias with EdgeValue. This gives six configurations and isolates
the addition of EdgeValue within each layerwise-bias feature setting. The
HLT-like campaign does not contain shared-bias-plus-EdgeValue models, so it
does not establish that layerwise bias is necessary for value messages.

The broader confirmation includes the baseline and capacity controls, every
single family, the all-family model, selected combinations, and the four
new architecture variants. A separately trained unary control supplies
endpoint information through particle embeddings without the explicit
pair-only descriptors. In total, 18 configurations receive three-seed final
test evaluation, covering 13 of the original screening configurations and
five additional configurations. The complete campaign therefore contains
26 distinct trained configurations, not 21 three-seed test configurations.
The screening seed is included in the three-seed comparison. Its checkpoint
is reused only when it matches the confirmation configuration; all other
seed/configuration instances are trained from random initialization.

Additional relation shuffling, wrong-event substitution, directional swaps,
and family or clustering-resolution removal are evaluated on comparison
validation. These are perturbations of trained models, not independently
retrained ablations. They test sensitivity to relation inputs and are not
interpreted as unbiased measures of each family's causal contribution.
Final-test metrics are reserved for reporting after checkpoint locking.

\subsection{HLT-like screening selects a compact relation set}
\label{sec:hlt-screening-results}

Table~\ref{tab:hlt-screening-results} reports every initial screening
configuration on comparison validation with one common training seed. PT+TRACK+REGION has the
highest observed accuracy, 74.6496\%, compared with 74.2008\% for the
standard baseline. It is the relation set selected by the prespecified
validation rule. TRACK and CHARGE give the largest single-family accuracy
estimates, at 74.5848\% and 74.5752\%, respectively. Their close values also
illustrate why the selection rule matters: CHARGE precedes TRACK under the
accuracy-window and cross-entropy ranking, although TRACK has the larger
accuracy point estimate.

% Single-seed validation; macro AUC averages all ten stored OVR AUCs.
\begin{table}[htbp]
\centering
\small
\setlength{\tabcolsep}{4pt}
\begin{tabularx}{\linewidth}{@{}Xrrr@{}}
\toprule
Configuration & Accuracy (\%) & $\Delta$ (pp) & Macro AUC \\
\midrule
PT+TRACK+REGION & 74.6496 & +0.4488 & 0.965217 \\ % data:screen:RPT_PT_TRACK_REGION
TRACK & 74.5848 & +0.3840 & 0.965276 \\ % data:screen:RPT_TRACK
CHARGE & 74.5752 & +0.3744 & 0.965199 \\ % data:screen:RPT_CHARGE
TRACK+PID & 74.5696 & +0.3688 & 0.965238 \\ % data:screen:RPT_TRACK_PID
PT & 74.5080 & +0.3072 & 0.964935 \\ % data:screen:RPT_PT
TRACK+PID+CHARGE & 74.4952 & +0.2944 & 0.965034 \\ % data:screen:RPT_TRACK_PID_CHARGE
All six families & 74.4792 & +0.2784 & 0.965115 \\ % data:screen:RPT_FULL_ALL
PT+REGION & 74.4576 & +0.2568 & 0.964819 \\ % data:screen:RPT_PT_REGION
PT+TRACK & 74.4504 & +0.2496 & 0.964921 \\ % data:screen:RPT_PT_TRACK
PID+CHARGE & 74.4360 & +0.2352 & 0.964788 \\ % data:screen:RPT_PID_CHARGE
PID & 74.4240 & +0.2232 & 0.965090 \\ % data:screen:RPT_PID
PT+TRACK+DENSITY & 74.4032 & +0.2024 & 0.964997 \\ % data:screen:RPT_PT_TRACK_DENSITY
TRACK+REGION & 74.3624 & +0.1616 & 0.964774 \\ % data:screen:RPT_TRACK_REGION
TRACK+CHARGE & 74.3616 & +0.1608 & 0.965037 \\ % data:screen:RPT_TRACK_CHARGE
DENSITY & 74.3592 & +0.1584 & 0.964795 \\ % data:screen:RPT_DENSITY
REGION & 74.3528 & +0.1520 & 0.964763 \\ % data:screen:RPT_REGION
PT+DENSITY & 74.3488 & +0.1480 & 0.964463 \\ % data:screen:RPT_PT_DENSITY
PT+TRACK+PID+CHARGE & 74.3360 & +0.1352 & 0.964833 \\ % data:screen:RPT_PT_TRACK_PID_CHARGE
Zero relations & 74.3176 & +0.1168 & 0.964554 \\ % data:screen:RPT_FULL_ZERO_REL
Wide base & 74.2784 & +0.0776 & 0.964656 \\ % data:screen:RPT_BASE_WIDE_MAX
Base & 74.2008 & --- & 0.964617 \\ % data:screen:RPT_BASE
\bottomrule
\end{tabularx}
\caption{All 21 HLT-like screening configurations on 125,000 comparison-
validation jets using one common training seed. All use shared pair bias without EdgeValue.
Rows are ordered by accuracy, not by the accuracy-window/cross-entropy
selection rule. Differences are relative to Base. Macro AUC is the
unweighted ten-class one-vs-rest mean on the 0--1 scale, including QCD.
All six families denotes PT+TRACK+PID+CHARGE+DENSITY+REGION; Wide base
and Zero relations are the capacity/input-shape controls.}
\label{tab:hlt-screening-results}
\end{table}


The selected configuration has macro AUC 0.965217, compared with 0.964617
for the baseline. TRACK has the largest screening macro AUC, 0.965276,
so the configuration with highest accuracy is not also the best on every
metric. Macro AUC and classwise rejection are reported descriptively;
they do not replace the original accuracy/cross-entropy selection rule.
Appendix~\ref{sec:hlt-screening-rejection-complete} reports rejection at
30\% and 50\% signal efficiency for every screening configuration and
every signal class on this same validation split.

Adding every relation family does not maximize validation accuracy: the
all-family configuration reaches 74.4792\%, below PT+TRACK+REGION. Conversely,
the relatively small standalone REGION improvement does not preclude its
presence in the selected combination. The selected combination exceeds the
best single-family accuracy by only 0.0648 percentage points in this screen.
We therefore use this experiment to choose a candidate representation, not
as conclusive evidence of feature complementarity or a universally optimal
relation set. The single-seed validation comparisons motivate the subsequent
multi-seed architecture study; they are not presented as final-test gains.

\subsection{HLT-like confirmation supports using relations in the value path}
\label{sec:hlt-architecture-results}

The six core architecture comparisons are shown in
Table~\ref{tab:hlt-architecture-results}. With standard pair features,
changing shared bias to layerwise bias produces a small mean test-accuracy
increase of 0.0175 percentage points. With the selected relations, the same
change increases accuracy by 0.1987 percentage points. Layerwise bias is
therefore not uniformly associated with a large gain across the two feature
settings.

% AUC is first averaged over classes per seed, then over seeds.
\begin{table}[htbp]
\centering
\footnotesize
\setlength{\tabcolsep}{4pt}
\begin{tabularx}{\linewidth}{@{}Xrrr@{}}
\toprule
Configuration & Accuracy (\%) & $\Delta$ (pp) & Macro AUC \\
\midrule
Standard: shared & $74.3744 \pm 0.0617$ & --- & $0.964941 \pm 0.000091$ \\ % data:hlt_arch:RPT_BASE
Standard: layerwise & $74.3919 \pm 0.0784$ & +0.0175 & $0.964931 \pm 0.000156$ \\ % data:hlt_arch:RPT_BASE_LAYERWISE
Standard: layerwise + EV & $74.7008 \pm 0.0929$ & +0.3264 & $0.965549 \pm 0.000285$ \\ % data:hlt_arch:RPT_BASE_EDGEVALUE
Selected: shared & $74.6366 \pm 0.0355$ & +0.2622 & $0.965375 \pm 0.000161$ \\ % data:hlt_arch:RPT_PT_TRACK_REGION
Selected: layerwise & $74.8353 \pm 0.0856$ & +0.4609 & $0.965846 \pm 0.000216$ \\ % data:hlt_arch:RPT_SELECTED_LAYERWISE
Selected: layerwise + EV & $75.0987 \pm 0.0438$ & +0.7243 & $0.966508 \pm 0.000126$ \\ % data:hlt_arch:RPT_SELECTED_EDGEVALUE
\bottomrule
\end{tabularx}
\caption{Six HLT-like architecture comparisons on the 500,000-jet final test.
Accuracy and macro AUC are three-seed means with sample standard deviations;
AUC uses the 0--1 scale. Selected features are standard features plus PT,
TRACK, and REGION. EV denotes EdgeValue; rows without EV use bias only.
Differences are relative to the standard shared-bias baseline.}
\label{tab:hlt-architecture-results}
\end{table}


Adding EdgeValue to an already-layerwise model increases mean accuracy by
0.3089 percentage points with standard features and 0.2633 percentage points
with selected features. The combined selected-relation, layerwise-bias,
EdgeValue model achieves a mean accuracy of 75.0987\% (seed SD 0.0438
percentage points), compared with 74.3744\% (SD 0.0617) for the baseline,
a mean improvement of 0.7243 percentage points. Its matched improvements are
positive at every seed, ranging from 0.6432 to 0.8060 percentage points.
These comparisons support the value-path modification within the tested
layerwise architectures; they do not test EdgeValue with shared bias.

Macro AUC increases from $0.964941\pm0.000091$ for the baseline to
$0.966508\pm0.000126$ for selected EdgeValue. The baseline layerwise-only
variant instead has macro AUC $0.964931\pm0.000156$, slightly below the
baseline mean despite its small increase in accuracy. The added AUC
comparison therefore also argues against treating the effect of layerwise
bias as uniformly favorable across metrics.

Appendix~\ref{sec:complete-hlt-results} reports all 18 three-seed final-test
configurations. The widened baseline and zero-relation controls improve
over the baseline by 0.0469 and 0.1087 percentage points, respectively,
while the trained selected-unary control improves by 0.0969 percentage
points. These particular controls do not reproduce the combined model's
accuracy gain. They do not, however, exclude all explanations involving
additional parameters or computation, because the architectural variants
are not all matched in both respects. The full results retain the relation
variants and controls even when their gains are small.

\subsection{HLT-like background rejection at fixed signal efficiency}
\label{sec:hlt-rejection-results}

Tables~\ref{tab:hlt-rejection-30} and~\ref{tab:hlt-rejection-50} compare
all nine signal classes for the six core HLT-like architectures at 30\%
and 50\% signal efficiency. These use the three-seed final-test
measurements, with 50,000 QCD jets per checkpoint. Complete rejection
tables for all 18 final-test configurations are provided in
Appendix~\ref{sec:hlt-test-rejection-complete}.

% Six core HLT architectures, original final test; all nine signals.
\begin{table}[htbp]
\centering
\footnotesize
\setlength{\tabcolsep}{2.5pt}
\begin{tabularx}{\linewidth}{@{}Xrrrrrrr@{}}
\toprule
Signal & B & BL & BE & S & SL & SE & SE/B change \\
\midrule
$H\to4q$ & 1062.72 & 1207.12 & 1308.06 & 1251.53 & 1293.30 & 1365.24 & +28.47\% \\ % data:rejection_hlt_core_30:H4q
$H\to b\bar b$ & 13055.56 & 11269.84 & 13888.89 & 9880.95 & 8492.06 & 11111.11 & -14.89\% \\ % data:rejection_hlt_core_30:Hbb
$H\to c\bar c$ & 1482.20 & 1608.55 & 1631.94 & 1525.72 & 1546.72 & 1924.98 & +29.87\% \\ % data:rejection_hlt_core_30:Hcc
$H\to gg$ & 142.77 & 144.69 & 145.28 & 153.34 & 153.73 & 154.08 & +7.92\% \\ % data:rejection_hlt_core_30:Hgg
$H\to q\bar q\ell\nu$ & 50000.00 & 50000.00 & 50000.00 & 50000.00 & SAT & SAT & --- \\ % data:rejection_hlt_core_30:Hqql
$t\to b\ell\nu$ & SAT & SAT & SAT & SAT & SAT & SAT & --- \\ % data:rejection_hlt_core_30:Tbl
$t\to bq\bar q'$ & 33333.33 & 24166.67 & 19444.44 & 41666.67 & 41666.67 & 22222.22 & -33.33\% \\ % data:rejection_hlt_core_30:Tbqq
$W\to q\bar q'$ & 519.04 & 534.83 & 599.25 & 520.63 & 509.29 & 594.67 & +14.57\% \\ % data:rejection_hlt_core_30:Wqq
$Z\to q\bar q$ & 537.68 & 479.27 & 470.39 & 489.08 & 460.81 & 522.85 & -2.76\% \\ % data:rejection_hlt_core_30:Zqq
\bottomrule
\end{tabularx}
\caption{HLT-like final-test QCD rejection at 30\% signal efficiency,
with 50,000 QCD jets and three seeds. B/BL/BE use standard features with
shared bias, layerwise bias, or layerwise bias plus EdgeValue; S/SL/SE
use the corresponding selected-feature configurations. Entries average
per-seed rejections. SAT means at least one seed passes zero QCD jets,
so no finite three-seed mean is reported. Changes use unrounded means.}
\label{tab:hlt-rejection-30}
\end{table}

% Six core HLT architectures, original final test; all nine signals.
\begin{table}[htbp]
\centering
\footnotesize
\setlength{\tabcolsep}{2.5pt}
\begin{tabularx}{\linewidth}{@{}Xrrrrrrr@{}}
\toprule
Signal & B & BL & BE & S & SL & SE & SE/B change \\
\midrule
$H\to4q$ & 310.97 & 300.40 & 305.89 & 308.71 & 299.45 & 311.95 & +0.31\% \\ % data:rejection_hlt_core_50:H4q
$H\to b\bar b$ & 2468.10 & 2090.91 & 2615.16 & 2266.08 & 2282.61 & 2500.00 & +1.29\% \\ % data:rejection_hlt_core_50:Hbb
$H\to c\bar c$ & 365.13 & 388.87 & 370.21 & 382.20 & 374.66 & 392.13 & +7.39\% \\ % data:rejection_hlt_core_50:Hcc
$H\to gg$ & 52.46 & 52.22 & 53.68 & 53.48 & 52.81 & 54.46 & +3.80\% \\ % data:rejection_hlt_core_50:Hgg
$H\to q\bar q\ell\nu$ & 25000.00 & 25000.00 & 25000.00 & 33333.33 & 33333.33 & 33333.33 & +33.33\% \\ % data:rejection_hlt_core_50:Hqql
$t\to b\ell\nu$ & SAT & SAT & SAT & SAT & SAT & SAT & --- \\ % data:rejection_hlt_core_50:Tbl
$t\to bq\bar q'$ & 2230.01 & 1954.16 & 2832.24 & 2511.96 & 2483.90 & 2738.10 & +22.78\% \\ % data:rejection_hlt_core_50:Tbqq
$W\to q\bar q'$ & 115.56 & 113.88 & 114.17 & 114.18 & 118.28 & 120.14 & +3.96\% \\ % data:rejection_hlt_core_50:Wqq
$Z\to q\bar q$ & 106.01 & 108.96 & 106.41 & 109.82 & 109.25 & 114.69 & +8.19\% \\ % data:rejection_hlt_core_50:Zqq
\bottomrule
\end{tabularx}
\caption{HLT-like final-test QCD rejection at 50\% signal efficiency,
with 50,000 QCD jets and three seeds. B/BL/BE use standard features with
shared bias, layerwise bias, or layerwise bias plus EdgeValue; S/SL/SE
use the corresponding selected-feature configurations. Entries average
per-seed rejections. SAT means at least one seed passes zero QCD jets,
so no finite three-seed mean is reported. Changes use unrounded means.}
\label{tab:hlt-rejection-50}
\end{table}


The selected EdgeValue model's higher accuracy and macro AUC do not imply
uniformly higher rejection. At 30\% signal efficiency it increases mean
rejection over baseline by 28.47\% for $H\to4q$ and 29.87\% for
$H\to c\bar c$, but decreases it by 14.89\% for $H\to b\bar b$,
33.33\% for $t\to bq\bar q'$, and 2.76\% for $Z\to q\bar q$.
At 50\% efficiency its mean rejection is higher at each operating point
where both means are finite. Saturated entries remain explicit, and
low-background-count operating points should not be interpreted as precise
population-level estimates. The screening rejection tables have still
smaller background support, 12,500 QCD jets, and remain labeled as
single-seed validation measurements rather than test results.

\clearpage
\subsection{Full HLT-like final-test comparison}
\label{sec:complete-hlt-results}

Table~\ref{tab:hlt-complete-results} retains every configuration in the
HLT-like final-test report. All rows use the same 500,000 test identities
and three training seeds. The 21 initial validation-screening rows are
reported separately in Table~\ref{tab:hlt-screening-results}. Together the
two tables cover all 26 distinct trained HLT-like configurations; they
should not be combined into a single ranking across validation and test.

% AUC is first averaged over classes per seed, then over seeds.
\begin{table}[htbp]
\centering
\footnotesize
\setlength{\tabcolsep}{4pt}
\begin{tabularx}{\linewidth}{@{}Xrrr@{}}
\toprule
Configuration & Accuracy (\%) & $\Delta$ (pp) & Macro AUC \\
\midrule
Selected + layerwise + EV & $75.0987 \pm 0.0438$ & +0.7243 & $0.966508 \pm 0.000126$ \\ % data:hlt_all:RPT_SELECTED_EDGEVALUE
Selected + layerwise & $74.8353 \pm 0.0856$ & +0.4609 & $0.965846 \pm 0.000216$ \\ % data:hlt_all:RPT_SELECTED_LAYERWISE
Base + layerwise + EV & $74.7008 \pm 0.0929$ & +0.3264 & $0.965549 \pm 0.000285$ \\ % data:hlt_all:RPT_BASE_EDGEVALUE
TRACK+PID & $74.6444 \pm 0.0232$ & +0.2700 & $0.965442 \pm 0.000057$ \\ % data:hlt_all:RPT_TRACK_PID
PT+TRACK+REGION & $74.6366 \pm 0.0355$ & +0.2622 & $0.965375 \pm 0.000161$ \\ % data:hlt_all:RPT_PT_TRACK_REGION
All six families & $74.6228 \pm 0.0489$ & +0.2484 & $0.965339 \pm 0.000130$ \\ % data:hlt_all:RPT_FULL_ALL
TRACK & $74.5694 \pm 0.1267$ & +0.1950 & $0.965403 \pm 0.000257$ \\ % data:hlt_all:RPT_TRACK
TRACK+CHARGE & $74.5590 \pm 0.1062$ & +0.1846 & $0.965358 \pm 0.000191$ \\ % data:hlt_all:RPT_TRACK_CHARGE
Zero relations & $74.4831 \pm 0.1077$ & +0.1087 & $0.964942 \pm 0.000267$ \\ % data:hlt_all:RPT_FULL_ZERO_REL
DENSITY & $74.4821 \pm 0.0270$ & +0.1077 & $0.964958 \pm 0.000018$ \\ % data:hlt_all:RPT_DENSITY
Selected unary & $74.4713 \pm 0.0716$ & +0.0969 & $0.964969 \pm 0.000185$ \\ % data:hlt_all:RPT_SELECTED_UNARY
CHARGE & $74.4705 \pm 0.0910$ & +0.0961 & $0.965114 \pm 0.000233$ \\ % data:hlt_all:RPT_CHARGE
REGION & $74.4575 \pm 0.0690$ & +0.0831 & $0.965046 \pm 0.000096$ \\ % data:hlt_all:RPT_REGION
PT & $74.4521 \pm 0.0499$ & +0.0777 & $0.965104 \pm 0.000087$ \\ % data:hlt_all:RPT_PT
Wide base & $74.4213 \pm 0.0881$ & +0.0469 & $0.964848 \pm 0.000130$ \\ % data:hlt_all:RPT_BASE_WIDE_MAX
PID & $74.4167 \pm 0.1103$ & +0.0423 & $0.965075 \pm 0.000276$ \\ % data:hlt_all:RPT_PID
Base + layerwise & $74.3919 \pm 0.0784$ & +0.0175 & $0.964931 \pm 0.000156$ \\ % data:hlt_all:RPT_BASE_LAYERWISE
Base & $74.3744 \pm 0.0617$ & --- & $0.964941 \pm 0.000091$ \\ % data:hlt_all:RPT_BASE
\bottomrule
\end{tabularx}
\caption{All 18 HLT-like final-test configurations on 500,000 jets, with three
matched training seeds. Accuracy and macro AUC are means with sample standard
deviations; macro AUC is on the 0--1 scale. Selected means PT+TRACK+REGION.
Selected unary is a separately trained endpoint-feature control.
These rows overlap with 13 initial screening configurations; they are
not 18 additional models.}
\label{tab:hlt-complete-results}
\end{table}


\clearpage
\subsection{Complete HLT-like screening rejection}
\label{sec:hlt-screening-rejection-complete}

Tables~\ref{tab:hlt-screen-rejection-1}--\ref{tab:hlt-screen-rejection-3}
give rejection at both 30\% and 50\% signal efficiency for all 21 initial
configurations. These are single-seed comparison-validation measurements on
125,000 jets, including 12,500 QCD jets. Signal classes are divided among
three panels for readability; every configuration appears in every panel.
All use shared bias without EdgeValue. SAT indicates zero observed QCD
background for the single checkpoint and does not establish infinite
population rejection.

% Complete HLT rejection coverage, grouped by signal class for readability.
\begin{table}[htbp]
\centering
\footnotesize
\setlength{\tabcolsep}{2.5pt}
\begin{tabularx}{\linewidth}{@{}Xrrrrrr@{}}
\toprule
Configuration & \multicolumn{2}{c}{$H\to4q$} & \multicolumn{2}{c}{$H\to b\bar b$} & \multicolumn{2}{c}{$H\to c\bar c$} \\
\cmidrule(lr){2-3}\cmidrule(lr){4-5}\cmidrule(lr){6-7}
 & 30\% & 50\% & 30\% & 50\% & 30\% & 50\% \\
\midrule
PT+TRACK+REGION & 892.86 & 290.70 & 4166.67 & 1388.89 & 1562.50 & 328.95 \\ % data:rejection_screen_g1:RPT_PT_TRACK_REGION
TRACK & 961.54 & 277.78 & 6250.00 & 1250.00 & 1785.71 & 284.09 \\ % data:rejection_screen_g1:RPT_TRACK
CHARGE & 892.86 & 357.14 & 6250.00 & 2500.00 & 1562.50 & 312.50 \\ % data:rejection_screen_g1:RPT_CHARGE
TRACK+PID & 694.44 & 250.00 & 4166.67 & 1785.71 & 1388.89 & 320.51 \\ % data:rejection_screen_g1:RPT_TRACK_PID
PT & 1136.36 & 265.96 & 4166.67 & 1388.89 & 961.54 & 277.78 \\ % data:rejection_screen_g1:RPT_PT
TRACK+PID+CHARGE & 1250.00 & 320.51 & 3125.00 & 1785.71 & 1136.36 & 284.09 \\ % data:rejection_screen_g1:RPT_TRACK_PID_CHARGE
All six families & 1136.36 & 297.62 & 2500.00 & 1388.89 & 1562.50 & 312.50 \\ % data:rejection_screen_g1:RPT_FULL_ALL
PT+REGION & 1136.36 & 284.09 & 4166.67 & 1250.00 & 1562.50 & 304.88 \\ % data:rejection_screen_g1:RPT_PT_REGION
PT+TRACK & 892.86 & 284.09 & 4166.67 & 1562.50 & 1388.89 & 265.96 \\ % data:rejection_screen_g1:RPT_PT_TRACK
PID+CHARGE & 1136.36 & 260.42 & 4166.67 & 1250.00 & 1250.00 & 312.50 \\ % data:rejection_screen_g1:RPT_PID_CHARGE
PID & 735.29 & 297.62 & 6250.00 & 1562.50 & 2083.33 & 337.84 \\ % data:rejection_screen_g1:RPT_PID
PT+TRACK+DENSITY & 833.33 & 250.00 & 6250.00 & 1785.71 & 1562.50 & 337.84 \\ % data:rejection_screen_g1:RPT_PT_TRACK_DENSITY
TRACK+REGION & 781.25 & 304.88 & 4166.67 & 1388.89 & 1136.36 & 297.62 \\ % data:rejection_screen_g1:RPT_TRACK_REGION
TRACK+CHARGE & 892.86 & 290.70 & 4166.67 & 1136.36 & 1250.00 & 337.84 \\ % data:rejection_screen_g1:RPT_TRACK_CHARGE
DENSITY & 961.54 & 320.51 & 6250.00 & 2083.33 & 1388.89 & 297.62 \\ % data:rejection_screen_g1:RPT_DENSITY
REGION & 1041.67 & 297.62 & 6250.00 & 2083.33 & 1388.89 & 357.14 \\ % data:rejection_screen_g1:RPT_REGION
PT+DENSITY & 833.33 & 290.70 & 4166.67 & 892.86 & 1250.00 & 312.50 \\ % data:rejection_screen_g1:RPT_PT_DENSITY
PT+TRACK+PID+CHARGE & 892.86 & 277.78 & 6250.00 & 1388.89 & 1562.50 & 337.84 \\ % data:rejection_screen_g1:RPT_PT_TRACK_PID_CHARGE
Zero relations & 1136.36 & 260.42 & 6250.00 & 1785.71 & 1388.89 & 320.51 \\ % data:rejection_screen_g1:RPT_FULL_ZERO_REL
Wide base & 892.86 & 260.42 & 2500.00 & 1136.36 & 1250.00 & 367.65 \\ % data:rejection_screen_g1:RPT_BASE_WIDE_MAX
Base & 961.54 & 260.42 & 6250.00 & 2083.33 & 1785.71 & 320.51 \\ % data:rejection_screen_g1:RPT_BASE
\bottomrule
\end{tabularx}
\caption{Complete HLT-like screening validation rejection (panel 1 of 3),
on 125,000 jets including 12,500 QCD jets. These are single-seed validation measurements, not final-test results.
SAT means zero passing QCD for that checkpoint.
The 30\% and 50\% columns refer to signal efficiency, not background
efficiency. Configuration names match the corresponding accuracy table.}
\label{tab:hlt-screen-rejection-1}
\end{table}

% Complete HLT rejection coverage, grouped by signal class for readability.
\begin{table}[htbp]
\centering
\footnotesize
\setlength{\tabcolsep}{2.5pt}
\begin{tabularx}{\linewidth}{@{}Xrrrrrr@{}}
\toprule
Configuration & \multicolumn{2}{c}{$H\to gg$} & \multicolumn{2}{c}{$H\to q\bar q\ell\nu$} & \multicolumn{2}{c}{$t\to b\ell\nu$} \\
\cmidrule(lr){2-3}\cmidrule(lr){4-5}\cmidrule(lr){6-7}
 & 30\% & 50\% & 30\% & 50\% & 30\% & 50\% \\
\midrule
PT+TRACK+REGION & 130.21 & 43.55 & SAT & SAT & SAT & SAT \\ % data:rejection_screen_g2:RPT_PT_TRACK_REGION
TRACK & 126.26 & 45.79 & SAT & SAT & SAT & SAT \\ % data:rejection_screen_g2:RPT_TRACK
CHARGE & 128.87 & 45.45 & SAT & SAT & SAT & SAT \\ % data:rejection_screen_g2:RPT_CHARGE
TRACK+PID & 137.36 & 45.62 & SAT & SAT & SAT & SAT \\ % data:rejection_screen_g2:RPT_TRACK_PID
PT & 113.64 & 45.13 & SAT & SAT & SAT & SAT \\ % data:rejection_screen_g2:RPT_PT
TRACK+PID+CHARGE & 120.19 & 45.79 & SAT & SAT & SAT & SAT \\ % data:rejection_screen_g2:RPT_TRACK_PID_CHARGE
All six families & 131.58 & 47.53 & SAT & SAT & SAT & SAT \\ % data:rejection_screen_g2:RPT_FULL_ALL
PT+REGION & 127.55 & 45.13 & SAT & SAT & SAT & SAT \\ % data:rejection_screen_g2:RPT_PT_REGION
PT+TRACK & 123.76 & 44.80 & SAT & SAT & SAT & SAT \\ % data:rejection_screen_g2:RPT_PT_TRACK
PID+CHARGE & 130.21 & 44.33 & SAT & SAT & SAT & SAT \\ % data:rejection_screen_g2:RPT_PID_CHARGE
PID & 128.87 & 45.13 & SAT & SAT & SAT & SAT \\ % data:rejection_screen_g2:RPT_PID
PT+TRACK+DENSITY & 126.26 & 42.81 & SAT & SAT & SAT & SAT \\ % data:rejection_screen_g2:RPT_PT_TRACK_DENSITY
TRACK+REGION & 131.58 & 45.45 & SAT & SAT & SAT & SAT \\ % data:rejection_screen_g2:RPT_TRACK_REGION
TRACK+CHARGE & 120.19 & 44.01 & SAT & SAT & SAT & SAT \\ % data:rejection_screen_g2:RPT_TRACK_CHARGE
DENSITY & 126.26 & 46.13 & SAT & SAT & SAT & SAT \\ % data:rejection_screen_g2:RPT_DENSITY
REGION & 142.05 & 45.62 & SAT & SAT & SAT & SAT \\ % data:rejection_screen_g2:RPT_REGION
PT+DENSITY & 126.26 & 43.10 & SAT & SAT & SAT & SAT \\ % data:rejection_screen_g2:RPT_PT_DENSITY
PT+TRACK+PID+CHARGE & 115.74 & 44.80 & SAT & SAT & SAT & SAT \\ % data:rejection_screen_g2:RPT_PT_TRACK_PID_CHARGE
Zero relations & 123.76 & 42.81 & SAT & SAT & SAT & SAT \\ % data:rejection_screen_g2:RPT_FULL_ZERO_REL
Wide base & 117.92 & 44.80 & SAT & 12500.00 & SAT & SAT \\ % data:rejection_screen_g2:RPT_BASE_WIDE_MAX
Base & 132.98 & 45.45 & SAT & SAT & SAT & SAT \\ % data:rejection_screen_g2:RPT_BASE
\bottomrule
\end{tabularx}
\caption{Complete HLT-like screening validation rejection (panel 2 of 3),
on 125,000 jets including 12,500 QCD jets. These are single-seed validation measurements, not final-test results.
SAT means zero passing QCD for that checkpoint.
The 30\% and 50\% columns refer to signal efficiency, not background
efficiency. Configuration names match the corresponding accuracy table.}
\label{tab:hlt-screen-rejection-2}
\end{table}

% Complete HLT rejection coverage, grouped by signal class for readability.
\begin{table}[htbp]
\centering
\footnotesize
\setlength{\tabcolsep}{2.5pt}
\begin{tabularx}{\linewidth}{@{}Xrrrrrr@{}}
\toprule
Configuration & \multicolumn{2}{c}{$t\to bq\bar q'$} & \multicolumn{2}{c}{$W\to q\bar q'$} & \multicolumn{2}{c}{$Z\to q\bar q$} \\
\cmidrule(lr){2-3}\cmidrule(lr){4-5}\cmidrule(lr){6-7}
 & 30\% & 50\% & 30\% & 50\% & 30\% & 50\% \\
\midrule
PT+TRACK+REGION & 12500.00 & 2083.33 & 480.77 & 121.36 & 378.79 & 100.81 \\ % data:rejection_screen_g3:RPT_PT_TRACK_REGION
TRACK & 12500.00 & 2500.00 & 446.43 & 115.74 & 462.96 & 105.04 \\ % data:rejection_screen_g3:RPT_TRACK
CHARGE & 6250.00 & 4166.67 & 446.43 & 126.26 & 367.65 & 91.24 \\ % data:rejection_screen_g3:RPT_CHARGE
TRACK+PID & 6250.00 & 2500.00 & 500.00 & 130.21 & 284.09 & 105.93 \\ % data:rejection_screen_g3:RPT_TRACK_PID
PT & 6250.00 & 3125.00 & 462.96 & 116.82 & 416.67 & 105.04 \\ % data:rejection_screen_g3:RPT_PT
TRACK+PID+CHARGE & 6250.00 & 3125.00 & 543.48 & 131.58 & 367.65 & 103.31 \\ % data:rejection_screen_g3:RPT_TRACK_PID_CHARGE
All six families & 12500.00 & 1785.71 & 520.83 & 117.92 & 390.62 & 104.17 \\ % data:rejection_screen_g3:RPT_FULL_ALL
PT+REGION & 6250.00 & 2083.33 & 500.00 & 117.92 & 367.65 & 96.90 \\ % data:rejection_screen_g3:RPT_PT_REGION
PT+TRACK & 6250.00 & 4166.67 & 500.00 & 114.68 & 320.51 & 95.42 \\ % data:rejection_screen_g3:RPT_PT_TRACK
PID+CHARGE & 6250.00 & 2500.00 & 390.62 & 113.64 & 328.95 & 96.15 \\ % data:rejection_screen_g3:RPT_PID_CHARGE
PID & 6250.00 & 2083.33 & 446.43 & 122.55 & 378.79 & 96.15 \\ % data:rejection_screen_g3:RPT_PID
PT+TRACK+DENSITY & 12500.00 & 3125.00 & 568.18 & 128.87 & 357.14 & 107.76 \\ % data:rejection_screen_g3:RPT_PT_TRACK_DENSITY
TRACK+REGION & 6250.00 & 3125.00 & 378.79 & 119.05 & 290.70 & 97.66 \\ % data:rejection_screen_g3:RPT_TRACK_REGION
TRACK+CHARGE & 12500.00 & 2500.00 & 462.96 & 117.92 & 347.22 & 104.17 \\ % data:rejection_screen_g3:RPT_TRACK_CHARGE
DENSITY & 12500.00 & 2500.00 & 480.77 & 116.82 & 337.84 & 104.17 \\ % data:rejection_screen_g3:RPT_DENSITY
REGION & SAT & 4166.67 & 520.83 & 123.76 & 367.65 & 99.21 \\ % data:rejection_screen_g3:RPT_REGION
PT+DENSITY & 12500.00 & 2500.00 & 480.77 & 112.61 & 357.14 & 99.21 \\ % data:rejection_screen_g3:RPT_PT_DENSITY
PT+TRACK+PID+CHARGE & 12500.00 & 2500.00 & 480.77 & 126.26 & 403.23 & 96.90 \\ % data:rejection_screen_g3:RPT_PT_TRACK_PID_CHARGE
Zero relations & 12500.00 & 2083.33 & 543.48 & 115.74 & 337.84 & 100.00 \\ % data:rejection_screen_g3:RPT_FULL_ZERO_REL
Wide base & 12500.00 & 3125.00 & 480.77 & 127.55 & 357.14 & 101.63 \\ % data:rejection_screen_g3:RPT_BASE_WIDE_MAX
Base & 12500.00 & 2500.00 & 568.18 & 114.68 & 390.62 & 102.46 \\ % data:rejection_screen_g3:RPT_BASE
\bottomrule
\end{tabularx}
\caption{Complete HLT-like screening validation rejection (panel 3 of 3),
on 125,000 jets including 12,500 QCD jets. These are single-seed validation measurements, not final-test results.
SAT means zero passing QCD for that checkpoint.
The 30\% and 50\% columns refer to signal efficiency, not background
efficiency. Configuration names match the corresponding accuracy table.}
\label{tab:hlt-screen-rejection-3}
\end{table}


\clearpage
\subsection{Complete HLT-like final-test rejection}
\label{sec:hlt-test-rejection-complete}

Tables~\ref{tab:hlt-hlt-rejection-1}--\ref{tab:hlt-hlt-rejection-3} cover
all 18 three-seed final-test configurations and all nine signal classes.
Each checkpoint sees 500,000 test jets, including 50,000 QCD jets. Entries
are arithmetic means of separately measured seed-level rejections, not
rejections obtained by pooling counts or ensembling predictions. A SAT
entry means at least one seed passes no QCD jets; finite seeds alone are
not averaged to fill such cells. These test tables must not be merged with
the smaller single-seed validation sample above.

% Complete HLT rejection coverage, grouped by signal class for readability.
\begin{table}[htbp]
\centering
\footnotesize
\setlength{\tabcolsep}{2.5pt}
\begin{tabularx}{\linewidth}{@{}Xrrrrrr@{}}
\toprule
Configuration & \multicolumn{2}{c}{$H\to4q$} & \multicolumn{2}{c}{$H\to b\bar b$} & \multicolumn{2}{c}{$H\to c\bar c$} \\
\cmidrule(lr){2-3}\cmidrule(lr){4-5}\cmidrule(lr){6-7}
 & 30\% & 50\% & 30\% & 50\% & 30\% & 50\% \\
\midrule
Selected + layerwise + EV & 1365.24 & 311.95 & 11111.11 & 2500.00 & 1924.98 & 392.13 \\ % data:rejection_hlt_g1:RPT_SELECTED_EDGEVALUE
Selected + layerwise & 1293.30 & 299.45 & 8492.06 & 2282.61 & 1546.72 & 374.66 \\ % data:rejection_hlt_g1:RPT_SELECTED_LAYERWISE
Base + layerwise + EV & 1308.06 & 305.89 & 13888.89 & 2615.16 & 1631.94 & 370.21 \\ % data:rejection_hlt_g1:RPT_BASE_EDGEVALUE
TRACK+PID & 1155.97 & 294.37 & 10277.78 & 2460.32 & 1724.87 & 388.82 \\ % data:rejection_hlt_g1:RPT_TRACK_PID
PT+TRACK+REGION & 1251.53 & 308.71 & 9880.95 & 2266.08 & 1525.72 & 382.20 \\ % data:rejection_hlt_g1:RPT_PT_TRACK_REGION
All six families & 1433.84 & 296.52 & 11666.67 & 2477.15 & 1729.63 & 386.95 \\ % data:rejection_hlt_g1:RPT_FULL_ALL
TRACK & 1281.25 & 298.92 & 8399.47 & 2146.46 & 1600.60 & 389.20 \\ % data:rejection_hlt_g1:RPT_TRACK
TRACK+CHARGE & 1335.98 & 303.90 & 10833.33 & 2479.02 & 1563.52 & 373.86 \\ % data:rejection_hlt_g1:RPT_TRACK_CHARGE
Zero relations & 1196.36 & 287.71 & 11111.11 & 2110.39 & 1616.16 & 366.83 \\ % data:rejection_hlt_g1:RPT_FULL_ZERO_REL
DENSITY & 1127.95 & 300.93 & 15000.00 & 2245.67 & 1539.70 & 368.58 \\ % data:rejection_hlt_g1:RPT_DENSITY
Selected unary & 1177.78 & 305.50 & 11269.84 & 2252.42 & 1605.37 & 398.98 \\ % data:rejection_hlt_g1:RPT_SELECTED_UNARY
CHARGE & 1221.45 & 313.88 & 16666.67 & 2698.41 & 1563.52 & 358.37 \\ % data:rejection_hlt_g1:RPT_CHARGE
REGION & 1251.60 & 303.08 & 9444.44 & 2518.52 & 1600.60 & 402.01 \\ % data:rejection_hlt_g1:RPT_REGION
PT & 1221.45 & 310.02 & 11111.11 & 2275.86 & 1569.66 & 383.78 \\ % data:rejection_hlt_g1:RPT_PT
Wide base & 1248.25 & 290.12 & 19444.44 & 2348.76 & 1583.39 & 357.79 \\ % data:rejection_hlt_g1:RPT_BASE_WIDE_MAX
PID & 1302.97 & 301.54 & 11111.11 & 2257.58 & 1670.26 & 354.98 \\ % data:rejection_hlt_g1:RPT_PID
Base + layerwise & 1207.12 & 300.40 & 11269.84 & 2090.91 & 1608.55 & 388.87 \\ % data:rejection_hlt_g1:RPT_BASE_LAYERWISE
Base & 1062.72 & 310.97 & 13055.56 & 2468.10 & 1482.20 & 365.13 \\ % data:rejection_hlt_g1:RPT_BASE
\bottomrule
\end{tabularx}
\caption{Complete HLT-like final test rejection (panel 1 of 3),
on 500,000 jets including 50,000 QCD jets. Entries are three-seed means of per-seed rejections, not pooled counts.
SAT means at least one zero-background seed and hence no finite mean.
The 30\% and 50\% columns refer to signal efficiency, not background
efficiency. Configuration names match the corresponding accuracy table.}
\label{tab:hlt-hlt-rejection-1}
\end{table}

% Complete HLT rejection coverage, grouped by signal class for readability.
\begin{table}[htbp]
\centering
\footnotesize
\setlength{\tabcolsep}{2.5pt}
\begin{tabularx}{\linewidth}{@{}Xrrrrrr@{}}
\toprule
Configuration & \multicolumn{2}{c}{$H\to gg$} & \multicolumn{2}{c}{$H\to q\bar q\ell\nu$} & \multicolumn{2}{c}{$t\to b\ell\nu$} \\
\cmidrule(lr){2-3}\cmidrule(lr){4-5}\cmidrule(lr){6-7}
 & 30\% & 50\% & 30\% & 50\% & 30\% & 50\% \\
\midrule
Selected + layerwise + EV & 154.08 & 54.46 & SAT & 33333.33 & SAT & SAT \\ % data:rejection_hlt_g2:RPT_SELECTED_EDGEVALUE
Selected + layerwise & 153.73 & 52.81 & SAT & 33333.33 & SAT & SAT \\ % data:rejection_hlt_g2:RPT_SELECTED_LAYERWISE
Base + layerwise + EV & 145.28 & 53.68 & 50000.00 & 25000.00 & SAT & SAT \\ % data:rejection_hlt_g2:RPT_BASE_EDGEVALUE
TRACK+PID & 142.78 & 52.28 & SAT & 50000.00 & SAT & SAT \\ % data:rejection_hlt_g2:RPT_TRACK_PID
PT+TRACK+REGION & 153.34 & 53.48 & 50000.00 & 33333.33 & SAT & SAT \\ % data:rejection_hlt_g2:RPT_PT_TRACK_REGION
All six families & 145.93 & 52.41 & SAT & 19444.44 & SAT & SAT \\ % data:rejection_hlt_g2:RPT_FULL_ALL
TRACK & 150.47 & 52.22 & SAT & 25000.00 & SAT & SAT \\ % data:rejection_hlt_g2:RPT_TRACK
TRACK+CHARGE & 150.42 & 53.01 & SAT & 25000.00 & SAT & SAT \\ % data:rejection_hlt_g2:RPT_TRACK_CHARGE
Zero relations & 140.20 & 51.79 & SAT & 22222.22 & SAT & SAT \\ % data:rejection_hlt_g2:RPT_FULL_ZERO_REL
DENSITY & 145.37 & 52.23 & SAT & 33333.33 & SAT & SAT \\ % data:rejection_hlt_g2:RPT_DENSITY
Selected unary & 148.04 & 50.91 & 50000.00 & 25000.00 & SAT & SAT \\ % data:rejection_hlt_g2:RPT_SELECTED_UNARY
CHARGE & 151.67 & 52.62 & SAT & 41666.67 & SAT & SAT \\ % data:rejection_hlt_g2:RPT_CHARGE
REGION & 150.45 & 52.25 & 50000.00 & 33333.33 & SAT & SAT \\ % data:rejection_hlt_g2:RPT_REGION
PT & 146.45 & 51.95 & 50000.00 & 41666.67 & SAT & SAT \\ % data:rejection_hlt_g2:RPT_PT
Wide base & 144.28 & 50.19 & SAT & 25000.00 & SAT & SAT \\ % data:rejection_hlt_g2:RPT_BASE_WIDE_MAX
PID & 144.70 & 52.58 & SAT & 33333.33 & SAT & SAT \\ % data:rejection_hlt_g2:RPT_PID
Base + layerwise & 144.69 & 52.22 & 50000.00 & 25000.00 & SAT & SAT \\ % data:rejection_hlt_g2:RPT_BASE_LAYERWISE
Base & 142.77 & 52.46 & 50000.00 & 25000.00 & SAT & SAT \\ % data:rejection_hlt_g2:RPT_BASE
\bottomrule
\end{tabularx}
\caption{Complete HLT-like final test rejection (panel 2 of 3),
on 500,000 jets including 50,000 QCD jets. Entries are three-seed means of per-seed rejections, not pooled counts.
SAT means at least one zero-background seed and hence no finite mean.
The 30\% and 50\% columns refer to signal efficiency, not background
efficiency. Configuration names match the corresponding accuracy table.}
\label{tab:hlt-hlt-rejection-2}
\end{table}

% Complete HLT rejection coverage, grouped by signal class for readability.
\begin{table}[htbp]
\centering
\footnotesize
\setlength{\tabcolsep}{2.5pt}
\begin{tabularx}{\linewidth}{@{}Xrrrrrr@{}}
\toprule
Configuration & \multicolumn{2}{c}{$t\to bq\bar q'$} & \multicolumn{2}{c}{$W\to q\bar q'$} & \multicolumn{2}{c}{$Z\to q\bar q$} \\
\cmidrule(lr){2-3}\cmidrule(lr){4-5}\cmidrule(lr){6-7}
 & 30\% & 50\% & 30\% & 50\% & 30\% & 50\% \\
\midrule
Selected + layerwise + EV & 22222.22 & 2738.10 & 594.67 & 120.14 & 522.85 & 114.69 \\ % data:rejection_hlt_g3:RPT_SELECTED_EDGEVALUE
Selected + layerwise & 41666.67 & 2483.90 & 509.29 & 118.28 & 460.81 & 109.25 \\ % data:rejection_hlt_g3:RPT_SELECTED_LAYERWISE
Base + layerwise + EV & 19444.44 & 2832.24 & 599.25 & 114.17 & 470.39 & 106.41 \\ % data:rejection_hlt_g3:RPT_BASE_EDGEVALUE
TRACK+PID & 37500.00 & 1972.31 & 557.59 & 117.15 & 509.76 & 106.44 \\ % data:rejection_hlt_g3:RPT_TRACK_PID
PT+TRACK+REGION & 41666.67 & 2511.96 & 520.63 & 114.18 & 489.08 & 109.82 \\ % data:rejection_hlt_g3:RPT_PT_TRACK_REGION
All six families & 22222.22 & 2176.54 & 559.77 & 116.92 & 523.00 & 108.08 \\ % data:rejection_hlt_g3:RPT_FULL_ALL
TRACK & 29166.67 & 2154.76 & 519.63 & 117.40 & 478.49 & 108.63 \\ % data:rejection_hlt_g3:RPT_TRACK
TRACK+CHARGE & 37500.00 & 2399.47 & 495.46 & 116.43 & 519.21 & 107.50 \\ % data:rejection_hlt_g3:RPT_TRACK_CHARGE
Zero relations & 30555.56 & 2182.54 & 499.17 & 113.33 & 490.61 & 103.47 \\ % data:rejection_hlt_g3:RPT_FULL_ZERO_REL
DENSITY & 16666.67 & 2197.07 & 525.89 & 114.33 & 459.36 & 107.62 \\ % data:rejection_hlt_g3:RPT_DENSITY
Selected unary & 19444.44 & 1928.21 & 551.71 & 119.21 & 493.81 & 105.84 \\ % data:rejection_hlt_g3:RPT_SELECTED_UNARY
CHARGE & 29166.67 & 2432.97 & 502.27 & 113.99 & 490.41 & 106.77 \\ % data:rejection_hlt_g3:RPT_CHARGE
REGION & 22222.22 & 2257.58 & 535.25 & 113.16 & 495.13 & 109.41 \\ % data:rejection_hlt_g3:RPT_REGION
PT & SAT & 2028.22 & 533.95 & 116.33 & 504.24 & 108.03 \\ % data:rejection_hlt_g3:RPT_PT
Wide base & 15277.78 & 1829.81 & 496.77 & 114.17 & 494.43 & 108.06 \\ % data:rejection_hlt_g3:RPT_BASE_WIDE_MAX
PID & SAT & 2278.74 & 532.71 & 116.55 & 515.61 & 108.96 \\ % data:rejection_hlt_g3:RPT_PID
Base + layerwise & 24166.67 & 1954.16 & 534.83 & 113.88 & 479.27 & 108.96 \\ % data:rejection_hlt_g3:RPT_BASE_LAYERWISE
Base & 33333.33 & 2230.01 & 519.04 & 115.56 & 537.68 & 106.01 \\ % data:rejection_hlt_g3:RPT_BASE
\bottomrule
\end{tabularx}
\caption{Complete HLT-like final test rejection (panel 3 of 3),
on 500,000 jets including 50,000 QCD jets. Entries are three-seed means of per-seed rejections, not pooled counts.
SAT means at least one zero-background seed and hence no finite mean.
The 30\% and 50\% columns refer to signal efficiency, not background
efficiency. Configuration names match the corresponding accuracy table.}
\label{tab:hlt-hlt-rejection-3}
\end{table}

